\documentclass[12pt, reqno]{amsart}

\usepackage{amsmath, amssymb, amsthm, amsfonts}

% avoid `Too many math alphabets used in version normal` error
\newcommand\hmmax{0}
\newcommand\bmmax{0}

\usepackage[mathscr]{euscript} 
\usepackage{stmaryrd} % St Mary's Road symbols font --- some extra symbols

%\usepackage{fontspec} 
\usepackage[xcharter]{newtxmath}
%\setmainfont{XCharter}

\usepackage{graphics, stackrel}
\usepackage{graphicx}

\usepackage{verbatim}
\usepackage{natbib}
\usepackage{enumitem}

%font
%\usepackage{lmodern}
%\usepackage[T1]{fontenc}
%\usepackage{fontspec} 
%\usepackage[xcharter]{newtxmath}
%\setmainfont{XCharter}
%\usepackage{mathpazo}
%\usepackage{tgpagella}

%subfloats / figures
\usepackage{caption}
\usepackage{subcaption}

% For pandas latex tables
\usepackage{booktabs}


\usepackage{fancyvrb}
\usepackage[dvipsnames,svgnames,table]{xcolor}
\usepackage{mdwlist}

\usepackage[breaklinks=true, citecolor=Brown, colorlinks=true, linkcolor=blue]{hyperref}

% lists
\usepackage{enumitem}
\setlist[enumerate]{itemsep=2pt,topsep=3pt}
\setlist[itemize]{itemsep=2pt,topsep=3pt}
\setlist[enumerate,1]{label=(\roman*)}

\usepackage{mathrsfs}  % caligraphic
%\usepackage{stix} 
\usepackage{bbm}
\usepackage{bm}        % bold symbols


%% page layout
\usepackage[left=1.25in, right=1.25in, top=1.0in, bottom=1.15in, includehead, includefoot]{geometry}

% nice inequalities
\renewcommand{\leq}{\leqslant}
\renewcommand{\geq}{\geqslant}

% inner product
\providecommand{\inner}[1]{\left\langle{#1}\right\rangle}
\providecommand{\innerp}[1]{\left\langle{#1}\right\rangle_\pi}

% lists
\usepackage{enumitem}
\setlist[enumerate]{itemsep=2pt,topsep=3pt}
\setlist[itemize]{itemsep=2pt,topsep=3pt}
\setlist[enumerate,1]{label={\upshape (\roman*)}}


\usepackage[ruled, linesnumbered]{algorithm2e}

%extra spacing
\renewcommand{\baselinestretch}{1.25}

%horizonal line
\newcommand{\HRule}{\rule{\linewidth}{0.3mm}}

% skip a line between paragraphs, no indentation
%\setlength{\parskip}{1.5ex plus0.5ex minus0.5ex}
%\setlength{\parindent}{0pt}

% footnote without a maker (blfootnote)
\newcommand\blfootnote[1]{%
  \begingroup
  \renewcommand\thefootnote{}\footnote{#1}%
  \addtocounter{footnote}{-1}%
  \endgroup
}

\DeclareMathOperator{\fix}{fix}
\DeclareMathOperator{\Span}{span}
\DeclareMathOperator{\diag}{diag}
\DeclareMathOperator*{\argmin}{arg\,min}
\DeclareMathOperator*{\argmax}{arg\,max}

\DeclareMathOperator{\cl}{cl}
\DeclareMathOperator{\Int}{int}
%\DeclareMathOperator{\overset{\circ}}{int}
\DeclareMathOperator{\Prob}{Prob}
\DeclareMathOperator{\determinant}{det}
\DeclareMathOperator{\Var}{Var}
\DeclareMathOperator{\Cov}{Cov}
\DeclareMathOperator{\graph}{graph}

% mics short cuts and symbols
\newcommand{\st}{\ensuremath{\ \mathrm{s.t.}\ }}
\newcommand{\setntn}[2]{ \{ #1 : #2 \} }
\newcommand{\fore}{\therefore \quad}
\newcommand{\preqsd}{\preceq_{sd} }
\newcommand{\toas}{\stackrel {\textrm{ \scriptsize{a.s.} }} {\to} }
\newcommand{\tod}{\stackrel { d } {\to} }
\newcommand{\too}{\stackrel { o } {\to} }
\newcommand{\toweak}{\stackrel { w } {\to} }
\newcommand{\topr}{\stackrel { p } {\to} }
\newcommand{\disteq}{\stackrel { \mathscr D } {=} }
\newcommand{\eqdist}{\stackrel {\textrm{ \scriptsize{d} }} {=} }
\newcommand{\iidsim}{\stackrel {\textrm{ {\sc iid }}} {\sim} }
\newcommand{\1}{\mathbbm 1}
\newcommand{\la}{\langle}
\newcommand{\ra}{\rangle}
\newcommand{\dee}{\,{\rm d}}
\newcommand{\og}{{\mathbbm G}}
\newcommand{\ctimes}{\! \times \!}
\newcommand{\sint}{{\textstyle\int}}

\newcommand{\given}{\, | \,}
\newcommand{\A}{\forall}

% d for integrals
\newcommand*\diff{\mathop{}\!\mathrm{d}}
\newcommand*\e{\mathrm{e}}

% nice emptyset
\let\oldemptyset\emptyset
\let\emptyset\varnothing


%\renewcommand{\times}{\! \times \!}

\newcommand{\aA}{\mathcal A}
\newcommand{\cC}{\mathscr C}
\newcommand{\sS}{\mathcal S}
\newcommand{\bB}{\mathcal B}
\newcommand{\oO}{\mathcal O}
\newcommand{\gG}{\mathcal G}
\newcommand{\hH}{\mathcal H}
\newcommand{\kK}{\mathcal K}
\newcommand{\iI}{\mathcal I}
\newcommand{\eE}{\mathcal E}
\newcommand{\fF}{\mathscr F}
\newcommand{\qQ}{\mathcal Q}
\newcommand{\tT}{\mathcal T}
\newcommand{\xX}{\mathcal X}
\newcommand{\yY}{\mathcal Y}
\newcommand{\rR}{\mathcal R}
\newcommand{\zZ}{\mathcal Z}
\newcommand{\wW}{\mathcal W}
\newcommand{\uU}{\mathcal U}
\newcommand{\lL}{\mathcal L}
\newcommand{\mM}{\mathcal M}
\newcommand{\dD}{\mathcal D}
\newcommand{\pP}{\mathcal P}
\newcommand{\vV}{\mathcal V}

\newcommand{\Bsf}{\mathsf B}
\newcommand{\Hsf}{\mathsf H}
\newcommand{\Vsf}{\mathsf V}

\newcommand{\BB}{\mathbbm B}
\newcommand{\DD}{\mathbbm D}
\newcommand{\RR}{\mathbbm R}
\newcommand{\CC}{\mathbbm C}
\newcommand{\QQ}{\mathbbm Q}
\newcommand{\NN}{\mathbbm N}
\newcommand{\GG}{\mathbbm G}
\newcommand{\UU}{\mathbbm U}
\newcommand{\TT}{\mathbbm T}
\newcommand{\YY}{\mathbbm Y}
\newcommand{\ZZ}{\mathbbm Z}
\newcommand{\HH}{\mathbbm H}
\newcommand{\MM}{\mathbbm M}
\newcommand{\PP}{\mathbbm P}
\newcommand{\VV}{\mathbbm V}
\newcommand{\EE}{\mathbbm E}


\newcommand{\bH}{\mathbf H}
\newcommand{\bT}{\mathbf T}

\newcommand{\var}{\mathbbm V}

\newcommand{\Xsf}{\mathsf X}
\newcommand{\Msf}{\mathsf M}
\newcommand{\Asf}{\mathsf A}
\newcommand{\Gsf}{\mathsf G}
\newcommand{\II}{\mathsf I}
\newcommand{\WW}{\mathsf W}

\renewcommand{\phi}{\varphi}
\renewcommand{\epsilon}{\varepsilon}

\newcommand{\bP}{\mathbf P}
\newcommand{\bQ}{\mathbf Q}
\newcommand{\bE}{\mathbf E}
\newcommand{\bM}{\mathbf M}
\newcommand{\bX}{\mathbf X}
\newcommand{\bY}{\mathbf Y}

\theoremstyle{plain}
\newtheorem{theorem}{Theorem}[section]
\newtheorem{corollary}[theorem]{Corollary}
\newtheorem{lemma}[theorem]{Lemma}
\newtheorem{proposition}[theorem]{Proposition}


\theoremstyle{definition}
\newtheorem{definition}{Definition}[section]
\newtheorem{axiom}{Axiom}[section]
\newtheorem{example}{Example}[section]
\newtheorem{remark}{Remark}[section]
\newtheorem{notation}{Notation}[section]
\newtheorem{assumption}{Assumption}[section]
\newtheorem{condition}{Condition}[section]


%\DeclareTextFontCommand{\emph}{\bfseries}

%%%%%%%%%%%%%%%%%% end my preamble %%%%%%%%%%%%%%%%%%%%%%%%%%%%%%%%%%%


\newcommand{\navy}[1]{\textcolor{blue}{#1}}




\begin{document}


\title{Semigroup Notes}

\author{JS and JY}

\date{\today}

\maketitle


In what follows, 

\begin{itemize}
    \item $E$ denotes a real Banach space,
    \item $\lL(E)$ is the set of bounded linear operators on $E$,
    \item $E'$ represents the dual space (i.e., the set of bounded linear
        functionals on $E$),
    \item $\| \cdot \|$ denotes either the norm on $E$ or the operator norm on
        $\lL(E)$, depending on context, and
    \item $\cC$ is the set of continuous functions from $\RR_+$ to $E$.
\end{itemize}

\section{Marcov Process}
\subsection{Transition probability}
Let $(X,\Sigma)$ be a measurable space and $(\Omega,\mathcal{F},\mathbb{P})$ be a probability space. A function $\pP\colon X\times \Sigma\rightarrow [0,1]$ is a \emph{transition kernel} if
\begin{enumerate} 
	\item\label{def: tk1} for all $B\in\Sigma$, $\pP(\cdot ,B)\colon X\rightarrow [0,1]$ is measurable w.r.t. $\Sigma.$
	\item \label{def: tk2}for all $x\in X$, $\pP(x,\cdot )\colon \Sigma\rightarrow [0,1]$ is a measure on $\Sigma.$
\end{enumerate}
If $\pP(x,X)=1$ for all $x\in X$, $\PP$ is a \emph{transition probability.} 

 We might need to consider the dynamics between different spaces so this definition is extended as follows. Let $(X_0,\Sigma_0)$ and  $(X_1,\Sigma_1)$ be two measurable spaces. Transition now is from $X_0$ to $X_1$ and so a function $\pP\colon X_0\times \Sigma_1\rightarrow [0,1]$ is a \emph{transition probability from $(X_0,\Sigma_0)$ to $(X_1,\Sigma_1)$ } if \ref{def: tk1} and \ref{def: tk2} hold for $\pP$ with the corresponding replacements.
 
 The following proposition shows the 'Kappa' can be constructed with a transition probability and a uniform distribution over $[0,1].$ ?
 
 \begin{proposition}
 	Let $\pP$ be a transition probability from a \emph{mesurable space} $(X_0,\Sigma_0)$ to a \emph{Boral space} $(X_1,\Sigma_1).$ There exists a measurable function $\kappa\colon X_0\times [0,1] \rightarrow X_1$ such that $\kappa(x,v)$ has the same distribution $\pP(x,\cdot)$ for every $x\in X_0$ if $v$ is uniformly distributed on $[0,1].$
 \end{proposition}
\begin{proof}
Since $(X_1,\Sigma_1)$ is a Boral space, there exist measureble bijection $\phi\colon X_1\rightarrow[0,1] .$ We need to define $\kappa$ $\st$ 
\begin{equation}\label{eq:kappa}
   \Prob(\kappa(x,v)\in B)=\pP(x,B)~~~\text{for all}~~x\in X_0,~ B\in \Sigma_1
\end{equation}
First let $F_x(t)=\pP(x,\phi^{-1}([0,t]))$ with $t\in [0,1]$ and so $t\mapsto F_x(t)$ is right continuous (same reason as CDF is r.c.) with $ F_x(1)=1.$ The generalized inverse of $F$ is
\begin{equation*}
   F^{\leftarrow}_x(q) =\inf\setntn{t\geq 0}{F_x(t)\geq q}, ~~q\geq0
\end{equation*}
Now we define 
\begin{equation*}
 \kappa(x,q)=\phi^{-1}(F^{\leftarrow}_x(q) )
\end{equation*}
Note that the probability distribution of a uniform random variable is the Lebesgue measure, the \eqref{eq:kappa} reduces to
\begin{equation}\label{eq:kappa2}
   Leb(q\in[0,1]: \kappa(x,q)\in \phi^{-1}([0,t]) )=\pP(x,\phi^{-1}([0,t]) )~~~\text{for all}~~x\in X_0,~ t\in [0,1]
\end{equation}
which is naturally true by definition of $\kappa$.
\end{proof}
\subsection{Some examples}
\subsubsection{Transition operators}
Let $\pP$ be a transition kernel on a measurable space $(X,\Sigma)$ and $B(X)$ be the Banach space of bounded measurable, real-valued functions on $X$ with supernorm
\begin{equation*}
	\|g\|=\sup_{x\in X}|g(x)|
\end{equation*}
\color{orange}(!! super norm is not the norm of $L^{\infty}$ )\color{black}

Given $g\in B(X)$, 
\begin{equation*}
	T g(x)=\int_{X}g(y) \diff\PP_x
\end{equation*}
where $\PP_x(B)=\pP(x,B)$ for all $B\in\Sigma.$ Then $T$ is a linear operator on $B(X).$ We define another linear operator $F$ on the set of all finite signed measures $\mathcal{U}(X)$ by 
\begin{equation*}
	F\mu(B)=\int_{X}\pP(x,B) \diff\mu
\end{equation*}
where $\mu\in\mathcal{U}^+(X)$ and it extends to all $ \mathcal{U}$ by linearity.

Now combining transition probability $\pP$ and linear operators $T,F$ together, we have for two $X$-valued r.v. $\psi_0$ and $\psi_1$ with distribution  $\mu_0$ and $\mu_1$. If $$\PP(\psi_1\in B|\psi_0=x)=\pP(x,B)$$ then we have
\begin{equation*}
	\mu_1=F\mu_0~~~\text{and}~~~\EE(g(\psi_1)|\psi_0=x))=\int_{X}g(y)\diff \PP_x=Tg(x)
\end{equation*}
for all $x\in X$ and $g\in B(X).$

\noindent\color{violet} Two linear operators \textbf{constructed from $\pP$} characterize the change of distributions of r.v. and property of conditional expectation on initial value.\color{black}


\subsubsection{Substochastic and stochastic operators}
Consider $L^1(X,\Sigma,\mu)$ and $L^{\infty}(X,\Sigma,\mu)$ with $\mu$ a $\sigma$-finite measure.  A linear operator $T\colon L^1(X,\Sigma,\mu)\rightarrow L^1(X,\Sigma,\mu)$ is \emph{substochastic} if $f\geq 0$ then $Tf\geq0$ for all $f\in L^1$ and if $\|Tf\|\leq \|f\|$ for all $f\in L^1$. The \emph{adjoint} operator $T^*\colon L^{\infty}(X,\Sigma,\mu)\rightarrow L^{\infty}(X,\Sigma,\mu)$ satisfies
\begin{equation*}
	\int_{X}Tf(x)g(x)\diff \mu=\int_{X}f(x)T^*g(x)\diff \mu
\end{equation*}
for all $f\in L^1$ and $g\in L^{\infty}.$ We now  connect $T$ and $T^*$ to transition kernel $\pP.$

\noindent\color{blue} Why are these two operators defined on different spaces? $L^{\infty}\subset L^1.$ Check which steps to make it well-defined \color{black}

\subsubsection{Integral stochastic operators}
%\subsubsection{Frobenius–Perron Operator}
\subsection{Discrete Marcov Processes}
\subsubsection{Marcov processes}
A family $(\xi_n)$ of $X$-valued r.v. has \emph{Markov property} if for all $n\geq0$
\begin{equation*}
    \PP(\xi_{n+1}\in B|\xi_0,\cdots,\xi_n)=\PP(\xi_{n=1}\in B|\xi_n)
\end{equation*}
\begin{proposition}
	A sequence $(\xi_n)$ has Marcov property if and only if for each $n$ and all $A\in \sigma(\xi_m: m\geq n)$ and  $F\in \sigma(\xi_m: m\leq n)$
	$$\PP(A\cup F|\xi_n)=\PP(A|\xi_n)\PP(F|\xi_n)$$
\end{proposition}
\begin{proof}
	A sequence of r.v. $(\xi_n)$ has Marcov propert. Then we have
	\begin{align*} 
	\PP(A\cup F|\xi_n)&=\PP(A\cup F|\xi_0,\cdots,\xi_n)~~~~~(\textcolor{orange}{Marcov~ property})\\
	                             &=\EE(\1_A\1_F|\xi_0,\cdots,\xi_n) \\
	                             &=\1_F\EE(\1_A|\xi_0,\cdots,\xi_n) ~~~~~(\textcolor{orange}{F\in  \sigma(\xi_m: m\leq n) })\\
	                             &=\EE(\1_F|\xi_0,\cdots,\xi_n)\EE(\1_A|\xi_0,\cdots,\xi_n)\\
	                             &=\PP(A|\xi_0,\cdots,\xi_n)\PP( F|\xi_0,\cdots,\xi_n)\\
	                             &=\PP(A|\xi_n)\PP( F|\xi_n)~~~~~(\textcolor{orange}{Marcov~ property})
	\end{align*}
	For the other direction, we let the equality holds. Let $A\in\sigma(\xi_n)$ and $F\in \sigma(\xi_m: m\leq n)$. The equality implies $\sigma(\xi_n)$ and $\sigma(\xi_m: m\leq n)$ are independent. Thus for $A\in \sigma(\xi_{n+1}),$
	\begin{equation*}
	\EE(\1_A|\xi_0,\cdots,\xi_n)=\EE(\EE(\1_A|\xi_n))|\xi_0,\cdots,\xi_n)=\EE(\1_A|\xi_n))
	\end{equation*}
\end{proof}
A stochastic process $(\xi_n)$ is a Marcov process with transition probability $\pP$ and initial distribution $\mu$ on $(X,\Sigma)$ if there is $(\Omega,\fF,\PP)$ $\st$ $\mu(B)=\PP(\xi_0\in B)$ and for all $B\in \Sigma,$
\begin{equation*}
      \PP(\xi_{n+1}\in B|\xi_0,\cdots,\xi_n)=\pP(\xi_n,B)
\end{equation*}
\textcolor{orange}{In this definition, transition probability $\pP$ gives the Marcov property and the initial distribution fixes the starting r.v.  }
%\begin{enumerate}
%	\item Go through appendix with conditional P and E on a sigma algebra
%	\item E and P meet when taking indicator functions
%	\item Visualize
%\end{enumerate}
% random mapping meets kappa again. dont really understand kappa....
\subsubsection{Canonical processes} % super familier name!!!
\newpage
\bibliographystyle{apalike}
\bibliography{sgs_bib}


\end{document}

